% !TeX program = pdflatex
% Compilation : pdflatex presentation_pfa.tex
% Présentation Beamer 16:9 -- Projet de Fin d'Année

\documentclass[10pt,aspectratio=169]{beamer}

\usepackage[utf8]{inputenc}
\usepackage[T1]{fontenc}
\usepackage[french]{babel}
\usepackage{lmodern}
\usepackage{booktabs}
\usepackage{amsmath}
\usepackage{tikz}
\usetikzlibrary{arrows.meta,positioning,shapes.geometric,fit,calc}

\definecolor{PfaBlue}{RGB}{31,95,151}
\definecolor{PfaDark}{RGB}{28,45,63}
\definecolor{PfaGreen}{RGB}{44,135,105}
\definecolor{PfaOrange}{RGB}{220,132,41}
\definecolor{PfaLight}{RGB}{237,244,249}

\usetheme{default}
\setbeamercolor{structure}{fg=PfaBlue}
\setbeamercolor{frametitle}{bg=PfaDark,fg=white}
\setbeamercolor{title}{fg=PfaDark}
\setbeamercolor{alerted text}{fg=PfaOrange}
\setbeamercolor{block title}{bg=PfaBlue,fg=white}
\setbeamercolor{block body}{bg=PfaLight,fg=black}
\setbeamertemplate{navigation symbols}{}
\setbeamertemplate{footline}{%
  \hfill\color{gray}\insertshorttitle\hspace{1em}--\hspace{1em}\insertframenumber/\inserttotalframenumber\hspace{1em}\vspace{0.6em}}
\setbeamertemplate{itemize items}[circle]
\setbeamersize{text margin left=0.8cm,text margin right=0.8cm}

\newcommand{\highlight}[1]{\textcolor{PfaBlue}{\textbf{#1}}}
\newcommand{\placeholder}[2]{%
  \begin{center}
    \fcolorbox{PfaBlue}{PfaLight}{\parbox[c][#2][c]{0.82\linewidth}{\centering\color{PfaDark}\textit{#1}}}
  \end{center}}

\title[Segmentation d'occupation du sol]{Développement d'un système intelligent de segmentation sémantique multi-classes d'images satellitaires}
\subtitle{Cartographie automatique de l'occupation du sol, post-traitement avancé et corrections manuelles persistantes}
\author{BOUMAZOUED Oumaima}
\institute{EMSI Rabat -- Ingénierie en Intelligence Artificielle et Data Science\\Entreprise d'accueil : Geomadev Morocco\\Encadrant : M. MOURCHID Ilyas}
\date{Année universitaire 2025--2026}

\begin{document}

% -----------------------------------------------------------------------------
\begin{frame}
  \titlepage
\end{frame}

\begin{frame}{Plan de la présentation}
  \tableofcontents
\end{frame}

% -----------------------------------------------------------------------------
\section{Contexte et objectifs}

\begin{frame}{Contexte}
  \begin{columns}[T]
    \column{0.56\textwidth}
      \begin{itemize}
        \item Les satellites et drones produisent de grands volumes d'images géographiques.
        \item Leur interprétation manuelle est lente, coûteuse et subjective.
        \item Les secteurs de l'urbanisme, de l'agriculture, de l'environnement et de la géomatique ont besoin de cartes fiables et rapides.
      \end{itemize}
      \vspace{0.25cm}
      \begin{block}{Idée centrale}
        Transformer une image RGB satellitaire ou aérienne en une carte d'occupation du sol exploitable.
      \end{block}
    \column{0.40\textwidth}
      \begin{tikzpicture}[node distance=0.5cm, every node/.style={align=center,font=\small}]
        \node[draw=PfaBlue,fill=PfaLight,rounded corners,minimum width=3.2cm,minimum height=0.8cm] (img) {Image satellite\\ou drone};
        \node[draw=PfaGreen,fill=green!8,rounded corners,below=of img,minimum width=3.2cm,minimum height=0.8cm] (ai) {Segmentation\\par IA};
        \node[draw=PfaOrange,fill=orange!8,rounded corners,below=of ai,minimum width=3.2cm,minimum height=0.8cm] (map) {Carte thématique\\et aide à la décision};
        \draw[-{Stealth[length=2mm]},thick,PfaBlue] (img) -- (ai);
        \draw[-{Stealth[length=2mm]},thick,PfaBlue] (ai) -- (map);
      \end{tikzpicture}
  \end{columns}
\end{frame}

\begin{frame}{Problématique et réponse proposée}
  \begin{block}{Problématique}
    Comment identifier de manière fiable les classes d'occupation du sol dans une image tout en permettant à l'utilisateur de corriger durablement les erreurs de prédiction ?
  \end{block}
  \vspace{0.25cm}
  \begin{columns}[T]
    \column{0.48\textwidth}
      \textbf{Difficultés observées}
      \begin{itemize}
        \item scènes hétérogènes et haute résolution ;
        \item ambiguïtés visuelles entre eau, forêt, agriculture et sol nu ;
        \item généralisation limitée à de nouvelles images ;
        \item modèle figé sans retour utilisateur.
      \end{itemize}
    \column{0.48\textwidth}
      \textbf{Réponse du projet}
      \begin{itemize}
        \item SegFormer-B0 comme modèle principal ;
        \item post-traitement fondé sur probabilités, texture, couleur et contexte spatial ;
        \item corrections manuelles persistantes, sans réentraînement ;
        \item plateforme web pour visualiser, corriger et exporter.
      \end{itemize}
  \end{columns}
\end{frame}

\begin{frame}{Objectifs}
  \begin{block}{Objectif principal}
    Concevoir une plateforme de cartographie automatique de l'occupation du sol à partir d'images satellitaires et aériennes.
  \end{block}
  \begin{enumerate}
    \item Prétraiter les images et assurer une inférence robuste sur des images de grande taille.
    \item Segmenter chaque pixel en classes d'occupation du sol avec SegFormer-B0.
    \item Réduire les confusions récurrentes Eau/Forêt/Agriculture/Sol nu, sans modifier les poids du modèle.
    \item Intégrer une correction manuelle persistante et réappliquée automatiquement.
    \item Proposer une interface web moderne, accessible aux utilisateurs SIG.
  \end{enumerate}
\end{frame}

% -----------------------------------------------------------------------------
\section{Données et modèle}

\begin{frame}{Jeu de données et classes}
  \begin{columns}[T]
    \column{0.48\textwidth}
      \begin{block}{Jeu de données LoveDA}
        \begin{itemize}
          \item Images RGB de télédétection haute résolution.
          \item Contextes urbains et ruraux.
          \item Masques annotés pour la segmentation sémantique.
        \end{itemize}
      \end{block}
      \vspace{0.2cm}
      \textbf{Prétraitement}\par
      Redimensionnement, normalisation ImageNet, transformations cohérentes image/masque et découpage en tuiles avec recouvrement pour les grandes images.
    \column{0.48\textwidth}
      \begin{block}{Classes affichées par la plateforme}
        \small
        \begin{tabular}{@{}ll@{\qquad}ll@{}}
          0 & Ignore & 4 & Water \\
          1 & Background & 5 & Barren \\
          2 & Building & 6 & Forest \\
          3 & Road & 7 & Agricultural \\
        \end{tabular}
      \end{block}
      \vspace{0.1cm}
      \alert{Les identifiants de classes sont conservés tout au long du pipeline.}
  \end{columns}
\end{frame}

\begin{frame}{Pourquoi SegFormer-B0 ?}
  \begin{columns}[T]
    \column{0.53\textwidth}
      \begin{itemize}
        \item Architecture Transformer hiérarchique, adaptée aux relations globales dans l'image.
        \item Fusion de caractéristiques à plusieurs échelles.
        \item Compromis intéressant entre qualité et coût d'inférence.
        \item Tête de segmentation adaptée aux 8 sorties du projet.
      \end{itemize}
      \vspace{0.25cm}
      \begin{alertblock}{Choix de production}
        SegFormer est le modèle par défaut. SegEarth-OV est disponible pour l'imagerie générale ; le mode hybride combine les deux approches.
      \end{alertblock}
    \column{0.43\textwidth}
      \begin{tikzpicture}[every node/.style={align=center,font=\small},node distance=0.35cm]
        \node[draw=PfaBlue,fill=PfaLight,rounded corners,minimum width=3.5cm] (a) {Image RGB};
        \node[draw=PfaBlue,fill=PfaLight,rounded corners,below=of a,minimum width=3.5cm] (b) {Encodeur Transformer\\hiérarchique};
        \node[draw=PfaBlue,fill=PfaLight,rounded corners,below=of b,minimum width=3.5cm] (c) {Décodeur MLP\\multi-échelles};
        \node[draw=PfaGreen,fill=green!8,rounded corners,below=of c,minimum width=3.5cm] (d) {Masque multi-classes};
        \draw[-{Stealth[length=2mm]},thick,PfaBlue] (a)--(b);
        \draw[-{Stealth[length=2mm]},thick,PfaBlue] (b)--(c);
        \draw[-{Stealth[length=2mm]},thick,PfaBlue] (c)--(d);
      \end{tikzpicture}
  \end{columns}
\end{frame}

% -----------------------------------------------------------------------------
\section{Architecture et pipeline}

\begin{frame}{Architecture globale de la solution}
  \centering
  \begin{tikzpicture}[
    node distance=0.58cm,
    box/.style={draw=PfaBlue,rounded corners,fill=PfaLight,minimum height=0.76cm,minimum width=3.1cm,align=center,font=\small},
    arrow/.style={-{Stealth[length=2mm]},thick,PfaBlue}]
    \node[box] (user) {Utilisateur};
    \node[box,below=of user] (ui) {Interface web\\React / TypeScript};
    \node[box,below=of ui] (api) {API Flask};
    \node[box,below left=0.85cm and 0.5cm of api] (model) {SegFormer / SegEarth-OV\\Mode hybride};
    \node[box,below right=0.85cm and 0.5cm of api] (post) {Post-traitement\\et corrections persistantes};
    \node[box,below=1.05cm of api] (out) {Carte de segmentation\\visualisation et export PNG};
    \draw[arrow] (user)--(ui);
    \draw[arrow] (ui)--(api);
    \draw[arrow] (api)--(model);
    \draw[arrow] (model)--(post);
    \draw[arrow] (post)--(out);
    \draw[arrow] (out.west) -| (ui.west);
  \end{tikzpicture}
\end{frame}

\begin{frame}{Pipeline de traitement}
  \centering
  \begin{tikzpicture}[
    node distance=0.18cm,
    step/.style={draw=PfaBlue,rounded corners,fill=PfaLight,minimum height=0.58cm,minimum width=10.8cm,align=left,font=\small},
    arrow/.style={-{Stealth[length=1.8mm]},thick,PfaBlue}]
    \node[step] (s1) {\textbf{1.} Chargement d'une image satellite ou drone};
    \node[step,below=of s1] (s2) {\textbf{2.} Prétraitement : conversion RGB, normalisation, redimensionnement et tuilage si nécessaire};
    \node[step,below=of s2] (s3) {\textbf{3.} Inférence : SegFormer, SegEarth-OV ou mode hybride};
    \node[step,below=of s3] (s4) {\textbf{4.} Raffinement : TTA, calibration, CRF/morphologie et levée des ambiguïtés};
    \node[step,below=of s4] (s5) {\textbf{5.} Réapplication des corrections déjà enregistrées pour cette image};
    \node[step,below=of s5] (s6) {\textbf{6.} Visualisation, correction interactive et export du résultat};
    \foreach \a/\b in {s1/s2,s2/s3,s3/s4,s4/s5,s5/s6} \draw[arrow] (\a)--(\b);
  \end{tikzpicture}
\end{frame}

\begin{frame}{Post-traitement avancé : réduire les confusions}
  \begin{columns}[T]
    \column{0.50\textwidth}
      \textbf{Problème ciblé}\par
      Certaines scènes végétalisées, agricoles ou de sol nu peuvent être prédites comme \highlight{Water}.
      \vspace{0.2cm}
      \begin{itemize}
        \item Calibration prudente des probabilités par classe.
        \item Comparaison Eau / Forêt / Agriculture / Barren.
        \item Analyse de texture locale et de couleur.
        \item Analyse de composantes : parcelles compactes contre rivières et lacs.
      \end{itemize}
    \column{0.46\textwidth}
      \begin{block}{Garde-fous}
        \begin{itemize}
          \item Les fortes prédictions d'eau sont conservées.
          \item Les rivières/canaux et plans d'eau lisses sont protégés.
          \item Les règles ne changent ni les IDs de classes ni les poids entraînés.
        \end{itemize}
      \end{block}
      \vspace{0.15cm}
      \alert{Objectif : améliorer les masques sans réentraîner le modèle.}
  \end{columns}
\end{frame}

\begin{frame}{Corrections manuelles persistantes}
  \begin{columns}[T]
    \column{0.52\textwidth}
      \begin{enumerate}
        \item L'utilisateur choisit la classe correcte.
        \item Il corrige la zone avec un \highlight{point}, un \highlight{pinceau} ou un \highlight{rectangle}.
        \item La correction est enregistrée avec l'empreinte SHA-256 du contenu de l'image.
        \item Lors d'une future analyse de la même image, elle est réappliquée automatiquement.
      \end{enumerate}
      \vspace{0.2cm}
      \begin{alertblock}{Bénéfice}
        Une boucle d'amélioration continue, sans fine-tuning et sans altérer les poids du modèle.
      \end{alertblock}
    \column{0.42\textwidth}
      \begin{tikzpicture}[node distance=0.45cm,every node/.style={align=center,font=\small}]
        \node[draw=PfaBlue,fill=PfaLight,rounded corners,minimum width=3.4cm] (a) {Prédiction};
        \node[draw=PfaOrange,fill=orange!10,rounded corners,below=of a,minimum width=3.4cm] (b) {Correction utilisateur};
        \node[draw=PfaGreen,fill=green!8,rounded corners,below=of b,minimum width=3.4cm] (c) {Stockage SHA-256};
        \node[draw=PfaBlue,fill=PfaLight,rounded corners,below=of c,minimum width=3.4cm] (d) {Réapplication automatique};
        \draw[-{Stealth[length=2mm]},thick,PfaBlue] (a)--(b);
        \draw[-{Stealth[length=2mm]},thick,PfaBlue] (b)--(c);
        \draw[-{Stealth[length=2mm]},thick,PfaBlue] (c)--(d);
      \end{tikzpicture}
  \end{columns}
\end{frame}

% -----------------------------------------------------------------------------
\section{Plateforme applicative}

\begin{frame}{Plateforme web développée}
  \begin{columns}[T]
    \column{0.48\textwidth}
      \textbf{Frontend React / TypeScript}
      \begin{itemize}
        \item chargement d'une image ou d'un dossier ;
        \item sélection du mode de segmentation ;
        \item carte de segmentation et superposition ;
        \item plein écran, mode sombre et interface bilingue FR/EN ;
        \item outils de correction accessibles dans la vue détaillée.
      \end{itemize}
    \column{0.48\textwidth}
      \textbf{Backend Flask / Python}
      \begin{itemize}
        \item chargement et cache des modèles ;
        \item inférence, tuilage et post-traitement ;
        \item API de prévisualisation, prédiction et correction ;
        \item stockage sûr des corrections ;
        \item génération de l'export PNG.
      \end{itemize}
  \end{columns}
  \vspace{0.25cm}
  \placeholder{Insérer ici une capture de l'interface : image originale + masque de segmentation + panneau de correction.}{2.0cm}
\end{frame}

\begin{frame}{Parcours utilisateur et export}
  \centering
  \begin{tikzpicture}[
    node distance=0.3cm,
    box/.style={draw=PfaBlue,rounded corners,fill=PfaLight,minimum height=0.85cm,minimum width=2.55cm,align=center,font=\small},
    arrow/.style={-{Stealth[length=2mm]},thick,PfaBlue}]
    \node[box] (a) {Choisir\\l'image};
    \node[box,right=of a] (b) {Choisir le\\mode d'analyse};
    \node[box,right=of b] (c) {Visualiser et\\corriger};
    \node[box,right=of c] (d) {Exporter\\le PNG};
    \draw[arrow] (a)--(b);
    \draw[arrow] (b)--(c);
    \draw[arrow] (c)--(d);
  \end{tikzpicture}
  \vspace{0.55cm}
  \begin{block}{Export implémenté}
    L'export ouvre le dialogue \emph{Enregistrer sous} du navigateur compatible (Chrome/Edge) : l'utilisateur choisit le dossier et le nom du fichier avant la sauvegarde de l'image PNG.
  \end{block}
  \small \textcolor{gray}{Les exports SIG vectoriels ou GeoTIFF constituent une évolution à finaliser dans l'implémentation actuelle.}
\end{frame}

% -----------------------------------------------------------------------------
\section{Évaluation et conclusion}

\begin{frame}{Évaluation : critères retenus}
  \begin{columns}[T]
    \column{0.48\textwidth}
      \begin{block}{Qualité de segmentation}
        \begin{itemize}
          \item Pixel Accuracy ;
          \item Precision, Recall et F1-score ;
          \item IoU par classe et mIoU ;
          \item Dice score.
        \end{itemize}
      \end{block}
    \column{0.48\textwidth}
      \begin{block}{Qualité opérationnelle}
        \begin{itemize}
          \item réduction des confusions Eau/Forêt/Agriculture/Barren ;
          \item robustesse CRF/TTA et tuilage ;
          \item temps d'inférence ;
          \item utilité des corrections persistantes.
        \end{itemize}
      \end{block}
  \end{columns}
  \vspace{0.25cm}
  \begin{alertblock}{À compléter avant la soutenance}
    Ajouter ici les valeurs mesurées : \texttt{Pixel Accuracy = ...}, \texttt{mIoU = ...}, temps moyen par image et comparaison avant/après post-traitement.
  \end{alertblock}
\end{frame}

\begin{frame}{Apports du projet}
  \begin{enumerate}
    \item Une chaîne complète, de l'image brute au masque exportable.
    \item Un modèle Transformer de production complété par des modes plus génériques.
    \item Un post-traitement explicable pour les classes visuellement ambiguës.
    \item Des corrections manuelles persistantes : l'outil mémorise l'expertise utilisateur.
    \item Une plateforme web orientée SIG, utilisable sans dépendre exclusivement d'un logiciel externe.
  \end{enumerate}
  \vspace{0.25cm}
  \begin{block}{Valeur ajoutée}
    Le projet dépasse une simple démonstration de modèle : il propose un outil de cartographie semi-supervisé, maintenable et évolutif.
  \end{block}
\end{frame}

\begin{frame}{Perspectives}
  \begin{columns}[T]
    \column{0.48\textwidth}
      \begin{itemize}
        \item Ajouter des exports GeoTIFF, GeoJSON et Shapefile réels.
        \item Exploiter des bandes multispectrales et indices comme NDVI/NDWI.
        \item Étendre les classes et les jeux de données.
      \end{itemize}
    \column{0.48\textwidth}
      \begin{itemize}
        \item Détection de changements à partir de séries temporelles.
        \item Optimiser l'inférence GPU et le déploiement cloud.
        \item Exploiter les corrections validées pour une future stratégie de fine-tuning contrôlé.
      \end{itemize}
  \end{columns}
\end{frame}

\begin{frame}{Conclusion}
  \begin{block}{Conclusion}
    La solution développée associe SegFormer, post-traitement avancé et correction manuelle persistante pour améliorer la cartographie automatique de l'occupation du sol à partir d'images satellitaires et aériennes.
  \end{block}
  \vfill
  \centering
  {\Large\color{PfaDark}\textbf{Merci pour votre attention}}\\[0.35cm]
  \large Questions ?
\end{frame}

\end{document}
